L'argument commence par la compatibilité des préfixes et des repères,
construit le code ternaire et détermine le drapeau au moyen de \(K\) et du
registre signé de \(\alpha\). L'incidence marquée sélectionne ensuite
l'origine réticulaire et le voisin de Leech. La dernière étape réalise l'algèbre
de vertex et ses automorphismes. À chaque étape, on précisera quelle structure
est conservée : combinaisons linéaires, supports, classes discriminantes, forme
bilinéaire ou produit gradué. Les explications de la fonction de \(K\),
de la normalisation et de l'origine de norme \(54\) accompagnent leurs
constructions respectives dans les sections \ref{exc:bandera} et \ref{exc:leech}.
